\chapter*{Anexo metodológico}
\addcontentsline{toc}{chapter}{Anexo metodológico}

\section*{Fuentes}

Criterios Generales de Política Económica 2025; Ley de Ingresos de la Federación
aprobada para 2024 y 2025 y su iniciativa; decreto de Presupuesto de Egresos de la
Federación; y analíticos presupuestarios del Presupuesto de Egresos aprobado de
2024 y 2025, en sus cuatro cortes, Gobierno Federal y entidades de control
directo, por programa y por función.

Cada cuadro de este documento se genera desde un archivo de datos con su registro
de procedencia: documento, ubicación y tier. Ninguna cifra del cuerpo del texto se
tecleó a mano.

\section*{Pruebas de cierre}

Las cifras de los Criterios Generales se leyeron de páginas renderizadas, porque
el documento se publicó sin capa de texto. La validación no fue una relectura sino
un conjunto de \textbf{cincuenta pruebas de identidad contable y de coherencia
entre anexos, todas las cuales pasan}. Entre ellas:

\begin{itemize}
\item Los requerimientos financieros igualan la suma de los recursos
extrapresupuestarios y el balance presupuestario.
\item El balance presupuestario iguala ingresos menos gasto neto pagado.
\item El gasto no programable iguala la suma de costo financiero, participaciones
y adeudos de ejercicios anteriores.
\item El balance primario iguala el balance presupuestario más el costo
financiero.
\item Las nueve razones a PIB de 2025 reproducen las publicadas usando el PIB
nominal declarado.
\item Los nueve conceptos que aparecen en dos anexos distintos coinciden en los
dos.
\item El gasto bruto del Gobierno Federal más el de las entidades, menos el
neteo del Anexo 1, iguala el gasto neto total del decreto.
\item El total de la Ley de Ingresos iguala el gasto neto total del decreto.
\item La diferencia entre los gastos obligatorios con y sin pensiones del Anexo 3
reproduce el agregado de pensiones y jubilaciones construido desde los analíticos,
al millón de pesos, en los dos ejercicios.
\end{itemize}

\section*{Convención de reparto de la descomposición}

El efecto denominador se reparte entre crecimiento e inflación como
$-d\,g/(1+n)$ y $-d\,\pi(1+g)/(1+n)$. Otra convención movería décimas entre las
dos variables sin cambiar su suma. La compensación por inflación pagada,
$+d\,\pi/(1+n)$, se contabiliza aparte porque viaja dentro del costo financiero y
de las adecuaciones a registros presupuestarios.

\textbf{Este documento no presenta el marco de manual de sostenibilidad}, primario
más tasa menos crecimiento, porque el paquete no publica un balance primario de
los requerimientos financieros y el del balance económico sesga al optimismo por
una o dos décimas.

\section*{Huecos declarados}

Población por afiliación para el gasto en salud por persona; población por entidad
federativa para el gasto federalizado por habitante; matrícula para el gasto
educativo por alumno; trayectoria de la tasa del derecho por la utilidad
compartida; y la exposición de motivos del proyecto de presupuesto, que no existe
para los ejercicios 2022 a 2026 porque está caída en el servidor de la Secretaría,
con las páginas índice enumerándola y los archivos devolviendo error.

\section*{Acrónimos}

CGPE, Criterios Generales de Política Económica. FASSA, Fondo de Aportaciones para
los Servicios de Salud. FONE, Fondo de Aportaciones para la Nómina Educativa y el
Gasto Operativo. IMSS, Instituto Mexicano del Seguro Social. ISSSTE, Instituto de
Seguridad y Servicios Sociales de los Trabajadores del Estado. LIF, Ley de
Ingresos de la Federación. PEF, Presupuesto de Egresos de la Federación. RFSPF,
requerimientos financieros del sector público federal, equivalente a RFSP. SHRFSPF,
saldo histórico de los requerimientos financieros del sector público federal,
equivalente a SHRFSP.


